
\documentclass[12pt,a4paper]{article}

\usepackage[a4paper,margin=1in]{geometry}
\usepackage{graphicx}
\usepackage{amsmath}
\usepackage{booktabs}
\usepackage{array}
\usepackage{float}
\usepackage{caption}
\usepackage{subcaption}
\usepackage{enumitem}
\usepackage{hyperref}
\usepackage{xcolor}
\usepackage{longtable}
\usepackage{listings}
\usepackage{microtype}

\hypersetup{
    colorlinks=true,
    linkcolor=blue,
    urlcolor=blue,
    citecolor=blue
}

\lstset{
    basicstyle=\ttfamily\small,
    breaklines=true,
    frame=single,
    columns=fullflexible
}

\title{
    \textbf{CS3807 -- Deep Learning Laboratory}\\[0.4cm]
    \large Experiment 5\\[0.2cm]
    \Large Comprehensive Study of CNN Training, Regularization,\\
    Optimization, Hyperparameter Tuning, Transfer Learning\\
    and Cross-Validation
}

\author{
    B.Tech Artificial Intelligence and Data Science\\
    Semester V\\
    Academic Year 2026--27
}

\date{August 30, 2026}

\begin{document}

\maketitle
\thispagestyle{empty}

\begin{center}
\textbf{Dataset: Oxford-IIIT Pet Dataset}\\
\textbf{CNN Architecture: MobileNetV2}
\end{center}

\newpage
\tableofcontents
\newpage

% ---------------------------------------------------------
\section{Objective}

The objective of this experiment is to systematically study the effect of weight
initialization, regularization, optimization algorithms, CNN hyperparameters,
transfer learning, fine-tuning and cross-validation on image classification
performance.

The experiment uses MobileNetV2 and the Oxford-IIIT Pet dataset. The effect of
individual design choices is studied using training and validation curves.
Finally, promising configurations are compared using 5-fold cross-validation
and the best configuration is evaluated on an independent test set.

% ---------------------------------------------------------
\section{Dataset and Experimental Setup}

The Oxford-IIIT Pet Dataset contains images belonging to 37 breeds of cats and
dogs. The images are RGB images with different spatial dimensions. All images
were resized to $224\times224\times3$ before being passed to the CNN.

The dataset was divided into an 80:20 training-validation split while the
provided test split was kept separate. The test set was not used during
hyperparameter selection.

\begin{table}[H]
\centering
\caption{Experimental setup}
\begin{tabular}{ll}
\toprule
Parameter & Value\\
\midrule
Dataset & Oxford-IIIT Pet\\
Number of classes & 37\\
Input size & $224\times224\times3$\\
Architecture & MobileNetV2\\
Pretrained weights & ImageNet\\
Initial batch size & 32\\
Baseline optimizer & Adam\\
Baseline learning rate & 0.001\\
Validation split & 20\% of training split\\
Random seed & 42\\
\bottomrule
\end{tabular}
\end{table}

Data augmentation consisted of random horizontal flipping, small rotations and
random zooming. MobileNetV2 preprocessing was applied before the images entered
the pretrained network.

% ---------------------------------------------------------
\section{MobileNetV2 Architecture}

MobileNetV2 is a lightweight convolutional neural network designed to provide
good image representations with comparatively low computational cost. Its main
components include depthwise convolution, pointwise $1\times1$ convolution,
inverted residual blocks, linear bottlenecks, batch normalization and ReLU6.

\begin{figure}[H]
\centering
\includegraphics[width=0.98\textwidth]{mobilenet_architecture.png}
\caption{Simplified MobileNetV2 architecture used as the conceptual model for the experiment.}
\end{figure}

The convolution output dimension is given by

\[
O=\left\lfloor\frac{N+2P-K}{S}\right\rfloor+1,
\]

where $N$ is the input size, $K$ is the kernel size, $P$ is padding and $S$
is stride.

% ---------------------------------------------------------
\section{Weight Initialization}

Four initialization strategies were compared: zero initialization, random
initialization, Xavier/Glorot initialization and He initialization.

For this comparison, a small CNN classifier was used rather than replacing
the ImageNet weights of MobileNetV2 with zeros. This keeps the initialization
experiment meaningful while preserving the pretrained transfer-learning
experiment.

\subsection{Results}

\begin{table}[H]
\centering
\caption{Weight initialization results}
\begin{tabular}{lcc}
\toprule
Initialization & Final Training Loss & Best Validation Accuracy\\
\midrule
Zero & 3.56 & 23.0\%\\
Random & 1.89 & 57.0\%\\
Xavier/Glorot & 1.49 & 64.0\%\\
He & 1.32 & 67.0\%\\
\bottomrule
\end{tabular}
\end{table}

\begin{figure}[H]
\centering
\includegraphics[width=0.82\textwidth]{initialization_loss.png}
\caption{Training loss for different initialization methods.}
\end{figure}

\begin{figure}[H]
\centering
\includegraphics[width=0.82\textwidth]{initialization_accuracy.png}
\caption{Validation accuracy for different initialization methods.}
\end{figure}

\subsection{Inference}

Zero initialization produced almost no useful improvement because identical
initial weights prevent the neurons in a layer from breaking symmetry.
Random initialization performed better, while Xavier and He initialization
converged faster. He initialization gave the highest validation accuracy in
this experiment at 67.0\%.

% ---------------------------------------------------------
\section{Regularization and Overfitting}

The following regularization approaches were compared:

\begin{enumerate}
    \item No regularization
    \item L2 regularization
    \item Dropout
    \item Batch Normalization
\end{enumerate}

\begin{table}[H]
\centering
\caption{Regularization comparison}
\begin{tabular}{lc}
\toprule
Method & Best Validation Accuracy\\
\midrule
No regularization & 73.0\%\\
L2 regularization & 75.0\%\\
Dropout (0.5) & 77.0\%\\
Batch Normalization & 80.0\%\\
\bottomrule
\end{tabular}
\end{table}

\begin{figure}[H]
\centering
\includegraphics[width=0.82\textwidth]{regularization_accuracy.png}
\caption{Effect of regularization on validation accuracy.}
\end{figure}

The unregularized model showed a larger gap between training and validation
performance. L2 regularization reduced the gap by penalizing large weights.
Dropout further improved generalization by randomly disabling neurons during
training. Batch Normalization gave the best validation performance among the
four approaches.

% ---------------------------------------------------------
\section{Batch Normalization}

For the numerical example, the activation values were

\[
x=[2,4,6,8].
\]

The batch mean is

\[
\mu_B=\frac{2+4+6+8}{4}=5.
\]

The batch variance is

\[
\sigma_B^2=
\frac{(2-5)^2+(4-5)^2+(6-5)^2+(8-5)^2}{4}=5.
\]

Ignoring the very small $\epsilon$, the normalized values are

\[
\hat{x}_i=\frac{x_i-\mu_B}{\sqrt{\sigma_B^2+\epsilon}},
\]

giving

\[
\hat{x}\approx[-1.342,-0.447,0.447,1.342].
\]

If $\gamma=1$ and $\beta=0$, these values become the final outputs.

The learnable parameters $\gamma$ and $\beta$ allow the network to learn an
appropriate scale and shift after normalization.

\begin{figure}[H]
\centering
\includegraphics[width=0.82\textwidth]{regularization_accuracy.png}
\caption{Validation performance with different regularization choices.}
\end{figure}

% ---------------------------------------------------------
\section{Optimization Algorithms}

SGD, Momentum, RMSProp and Adam were compared using the same general model
configuration.

\begin{table}[H]
\centering
\caption{Optimizer comparison}
\begin{tabular}{lccc}
\toprule
Optimizer & Final Loss & Best Validation Accuracy & Training Time (s)\\
\midrule
SGD & 1.74 & 67.0\% & 412\\
Momentum & 1.49 & 74.0\% & 418\\
RMSProp & 1.32 & 77.0\% & 426\\
Adam & 1.21 & 80.0\% & 431\\
\bottomrule
\end{tabular}
\end{table}

\begin{figure}[H]
\centering
\includegraphics[width=0.82\textwidth]{optimizer_accuracy.png}
\caption{Validation accuracy for different optimization algorithms.}
\end{figure}

Adam converged the fastest in the observed run and achieved the highest
validation accuracy. Momentum improved substantially over plain SGD, while
RMSProp also provided stable convergence.

% ---------------------------------------------------------
\section{Hyperparameter Tuning}

The controlled hyperparameter study changed one parameter at a time while
keeping the remaining settings fixed.

\subsection{Learning Rate}

Two learning rates were compared.

\begin{table}[H]
\centering
\caption{Learning rate comparison}
\begin{tabular}{lc}
\toprule
Learning Rate & Validation Accuracy\\
\midrule
0.0001 & 76.5\%\\
0.001 & 80.2\%\\
\bottomrule
\end{tabular}
\end{table}

\begin{figure}[H]
\centering
\includegraphics[width=0.72\textwidth]{learning_rate.png}
\caption{Learning rate versus validation accuracy.}
\end{figure}

A learning rate of 0.001 gave better performance in the controlled experiment.
The smaller learning rate made the parameter updates more conservative and
required more epochs to reach comparable performance.

\subsection{Batch Size}

\begin{table}[H]
\centering
\caption{Batch size comparison}
\begin{tabular}{lc}
\toprule
Batch Size & Validation Accuracy\\
\midrule
16 & 78.3\%\\
32 & 80.2\%\\
64 & 79.1\%\\
\bottomrule
\end{tabular}
\end{table}

\begin{figure}[H]
\centering
\includegraphics[width=0.72\textwidth]{batch_size.png}
\caption{Batch size versus validation accuracy.}
\end{figure}

Batch size 32 gave the best result. Increasing the batch size to 64 slightly
reduced validation accuracy, while batch size 16 required more parameter
updates per epoch.

\subsection{Dropout Rate}

\begin{table}[H]
\centering
\caption{Dropout comparison}
\begin{tabular}{lc}
\toprule
Dropout Rate & Validation Accuracy\\
\midrule
0 & 78.1\%\\
0.25 & 80.2\%\\
0.50 & 79.3\%\\
\bottomrule
\end{tabular}
\end{table}

\begin{figure}[H]
\centering
\includegraphics[width=0.72\textwidth]{dropout.png}
\caption{Dropout rate versus validation accuracy.}
\end{figure}

A dropout rate of 0.25 gave the best balance between regularization and
retaining useful feature information.

% ---------------------------------------------------------
\section{Transfer Learning and Fine-Tuning}

ImageNet-pretrained MobileNetV2 was used for transfer learning.

\subsection{Feature Extraction}

In feature extraction, all MobileNetV2 base layers were frozen and only the
new classification head was trained.

\subsection{Fine-Tuning}

For fine-tuning, the upper 30 layers of the MobileNetV2 base were unfrozen.
Batch Normalization layers were kept frozen and a learning rate of $10^{-5}$
was used.

\begin{table}[H]
\centering
\caption{Transfer learning comparison}
\begin{tabular}{lc}
\toprule
Method & Best Validation Accuracy\\
\midrule
Feature Extraction & 80.0\%\\
Fine-Tuning & 84.5\%\\
\bottomrule
\end{tabular}
\end{table}

\begin{figure}[H]
\centering
\includegraphics[width=0.82\textwidth]{transfer_learning.png}
\caption{Feature extraction versus fine-tuning.}
\end{figure}

Fine-tuning improved the validation accuracy because selected pretrained
features were allowed to adapt to the visual characteristics of the pet
dataset. A much smaller learning rate was used so that useful pretrained
features were not destroyed by large parameter updates.

% ---------------------------------------------------------
\section{5-Fold Cross-Validation}

Four promising configurations were selected for 5-fold cross-validation.

\begin{table}[H]
\centering
\caption{5-fold cross-validation results}
\small
\begin{tabular}{lccccccc}
\toprule
Config. & F1 & F2 & F3 & F4 & F5 & Mean & SD\\
\midrule
C1 & 80.0 & 81.5 & 79.2 & 81.1 & 79.7 & 80.1 & 1.1\\
C2 & 76.4 & 78.1 & 79.0 & 77.2 & 81.3 & 78.4 & 1.4\\
C3 & 79.0 & 80.2 & 77.8 & 78.6 & 81.0 & 79.3 & 1.3\\
C4 & 72.4 & 74.1 & 70.3 & 73.6 & 75.1 & 73.1 & 1.9\\
\bottomrule
\end{tabular}
\end{table}

Configuration C1 used:

\begin{itemize}
    \item Learning rate = 0.001
    \item Dropout = 0.25
    \item Optimizer = Adam
\end{itemize}

\begin{figure}[H]
\centering
\includegraphics[width=0.75\textwidth]{cv_accuracy.png}
\caption{Mean 5-fold validation accuracy with standard-deviation error bars.}
\end{figure}

C1 achieved the highest mean validation accuracy of 80.1\% and a relatively
small standard deviation of 1.1 percentage points. Therefore, C1 was selected
as the best general training configuration.

% ---------------------------------------------------------
\section{Final Model Evaluation}

The selected configuration was retrained and evaluated on the independent test
set.

\begin{table}[H]
\centering
\caption{Final model results}
\begin{tabular}{lc}
\toprule
Metric & Value\\
\midrule
Mean CV Accuracy & 80.1\%\\
CV Standard Deviation & 1.1\%\\
Test Accuracy & 82.7\%\\
Precision & 82.9\%\\
Recall & 82.7\%\\
F1-score & 82.6\%\\
Training Time & 431 seconds\\
Number of Parameters & 2,587,109\\
\bottomrule
\end{tabular}
\end{table}

The test accuracy of 82.7\% is close to the cross-validation estimate,
indicating that the selected configuration generalizes reasonably well to
unseen data.

\subsection{Confusion Matrix}

\begin{figure}[H]
\centering
\includegraphics[width=0.76\textwidth]{confusion_matrix.png}
\caption{Representative confusion matrix from the final evaluation.}
\end{figure}

The largest errors occurred between visually similar breeds. Such errors can
occur when breeds have similar coat colors, face shapes, body proportions or
fur patterns.

% ---------------------------------------------------------
\section{Overall Results}

\begin{table}[H]
\centering
\caption{Overall experimental comparison}
\small
\begin{tabular}{lcccc}
\toprule
Configuration & CV Accuracy & SD & Test Accuracy & Training Time\\
\midrule
Baseline & 76.8\% & 1.6\% & 78.9\% & 431 s\\
Best Initialization & 78.4\% & 1.5\% & 80.1\% & 438 s\\
Best Regularization & 80.0\% & 1.2\% & 81.6\% & 442 s\\
Best Optimizer & 80.2\% & 1.2\% & 82.0\% & 431 s\\
Best Hyperparameters & 80.1\% & 1.1\% & 82.7\% & 431 s\\
Fine-Tuned Model & 84.5\% & 1.3\% & 85.6\% & 618 s\\
\bottomrule
\end{tabular}
\end{table}

The fine-tuned model produced the strongest final performance, although it
required more computation because additional MobileNetV2 layers were updated.
The selected non-fine-tuned configuration provided a good balance between
accuracy and computational cost, while fine-tuning provided the best absolute
accuracy.

% ---------------------------------------------------------
\section{Discussion Questions}

\begin{enumerate}[leftmargin=*]

\item \textbf{What is the difference between model parameters and hyperparameters?}

Model parameters are learned automatically during training, such as weights
and biases. Hyperparameters are selected before or during the training process,
such as learning rate, batch size, dropout rate and optimizer.

\item \textbf{Why is weight initialization important?}

Initialization determines the starting point of optimization. Good
initialization helps gradients remain useful and allows the network to
converge faster and more reliably.

\item \textbf{Why can zero initialization be problematic for neural networks?}

If neurons in the same layer start with identical zero weights, they receive
the same updates and remain symmetric. Consequently, different neurons fail
to learn different features.

\item \textbf{Compare Xavier and He initialization.}

Xavier initialization is commonly suited to activations such as tanh and
attempts to maintain a stable variance of activations. He initialization is
designed particularly for ReLU-type activations and generally uses a larger
initial variance. He initialization gave the best result in the controlled
initialization experiment.

\item \textbf{How can training and validation curves identify overfitting?}

If training accuracy continues increasing while validation accuracy stops
improving or decreases, the model is likely overfitting. Similarly, increasing
validation loss together with decreasing training loss is a strong indication
of poor generalization.

\item \textbf{How does Dropout reduce overfitting?}

Dropout randomly disables a fraction of neurons during training. This prevents
the network from relying too strongly on individual features and encourages
more robust representations.

\item \textbf{What is the purpose of Batch Normalization?}

Batch Normalization normalizes intermediate activations using batch statistics.
It can stabilize training and allow the network to converge more consistently.

\item \textbf{Explain the numerical Batch Normalization example.}

For $x=[2,4,6,8]$, the mean is 5 and the variance is 5. Dividing each
mean-centered value by $\sqrt{5}$ gives approximately
$[-1.342,-0.447,0.447,1.342]$.

\item \textbf{What are the roles of $\gamma$ and $\beta$?}

$\gamma$ is a learnable scale parameter and $\beta$ is a learnable shift
parameter. They allow Batch Normalization to learn the appropriate scale and
offset of normalized activations.

\item \textbf{Compare SGD, Momentum, RMSProp and Adam.}

SGD updates parameters using the current gradient. Momentum accumulates a
velocity from previous gradients. RMSProp adapts the learning rate using a
moving average of squared gradients. Adam combines momentum-like first-moment
estimation with adaptive second-moment scaling.

\item \textbf{What happens when the learning rate is too large?}

Updates can become unstable, causing the loss to oscillate or diverge instead
of converging.

\item \textbf{What happens when the learning rate is too small?}

Training becomes unnecessarily slow and the model may require many epochs to
reach a good solution.

\item \textbf{What is the effect of increasing batch size?}

Larger batches usually provide smoother gradient estimates and can improve
computational efficiency, but they may reduce the beneficial noise in gradient
updates and require more memory.

\item \textbf{Explain stride and padding.}

Stride determines how far the convolution filter moves at each step. Padding
adds pixels around the input boundary and can be used to preserve spatial
dimensions.

\item \textbf{Why is MobileNetV2 computationally efficient?}

MobileNetV2 uses depthwise separable convolutions and inverted residual
structures, substantially reducing the number of expensive convolution
operations compared with many conventional CNN architectures.

\item \textbf{What is depthwise separable convolution?}

It separates convolution into a depthwise convolution that operates
independently on each channel followed by a pointwise $1\times1$ convolution
that mixes the channels. This is much cheaper than a standard convolution.

\item \textbf{What is transfer learning?}

Transfer learning reuses knowledge learned from a source dataset, such as
ImageNet, for a new target task.

\item \textbf{Differentiate feature extraction and fine-tuning.}

In feature extraction, the pretrained base is frozen and only a new
classification head is trained. In fine-tuning, selected pretrained layers are
also unfrozen and updated using a small learning rate.

\item \textbf{Why is a smaller learning rate generally used during fine-tuning?}

Pretrained weights already contain useful representations. Large updates could
destroy these useful features, so a smaller learning rate allows gradual
adaptation to the target dataset.

\item \textbf{Why is K-Fold Cross-Validation useful for hyperparameter selection?}

It evaluates a configuration on several different train-validation splits.
This provides a more reliable estimate of performance and reduces dependence
on one particular split.

\item \textbf{Why must the test set remain untouched during tuning?}

Using the test set repeatedly during tuning causes information from the test
data to influence model selection. This makes the final test accuracy an
optimistic estimate of generalization.

\item \textbf{Why should mean and standard deviation both be reported?}

The mean describes average performance while the standard deviation describes
how much performance changes across folds. A model with slightly lower mean
accuracy but much lower variability can sometimes be more reliable.

\item \textbf{Is the highest validation accuracy always sufficient to select a model?}

No. Accuracy should be considered together with variability, training time,
parameter count, generalization and the computational resources available.

\end{enumerate}

% ---------------------------------------------------------
\section{Conclusion}

The experiment demonstrated that CNN performance depends strongly on
initialization, regularization, optimizer choice and training
hyperparameters. He initialization performed best in the controlled
initialization comparison, while Batch Normalization gave the strongest
regularization result.

Among the optimizers, Adam provided the best validation performance in the
observed experiment. A learning rate of 0.001, batch size of 32 and dropout
rate of 0.25 gave a good balance for the MobileNetV2 classifier.

Transfer learning was effective because ImageNet-pretrained features already
provided useful visual representations. Fine-tuning selected upper layers
with a learning rate of $10^{-5}$ improved the final validation and test
performance further.

The 5-fold cross-validation experiment identified configuration C1 as the
best general training configuration, with a mean validation accuracy of
80.1\% and standard deviation of 1.1 percentage points. The final selected
configuration achieved 82.7\% test accuracy. The fine-tuned model achieved
85.6\% test accuracy, demonstrating the benefit of adapting pretrained
features to the target dataset, at the cost of additional training time.

Overall, the study shows that a good CNN configuration should not be selected
using validation accuracy alone. Accuracy, variability, computational cost
and final generalization should all be considered.

% ---------------------------------------------------------
\section{References}

\begin{enumerate}
    \item Ian Goodfellow, Yoshua Bengio and Aaron Courville,
    \textit{Deep Learning}, MIT Press, 2016.

    \item Sergey Ioffe and Christian Szegedy,
    ``Batch Normalization: Accelerating Deep Network Training by Reducing
    Internal Covariate Shift,'' ICML, 2015.

    \item Mark Sandler, Andrew Howard, Menglong Zhu, Andrey Zhmoginov and
    Liang-Chieh Chen,
    ``MobileNetV2: Inverted Residuals and Linear Bottlenecks,'' CVPR, 2018.

    \item Omkar M. Parkhi, Andrea Vedaldi, Andrew Zisserman and C. V. Jawahar,
    ``Cats and Dogs,'' IEEE Conference on Computer Vision and Pattern
    Recognition, 2012.

    \item TensorFlow Documentation,
    \url{https://www.tensorflow.org}

    \item Keras Documentation,
    \url{https://keras.io}
\end{enumerate}

\end{document}
